\VerdierPhysicalPageAnchor{0059}%
\VerdierSemanticAnchor{ch1-1-5-13}
\noindent\textbf{1.5.13.}\quad
Les $\mathbb Z^n$-catégories que nous utiliserons dans ce travail seront
des catégories strictes. Aussi ne tenterons-nous pas d'aborder le
problème signalé en (1.5.12) dans le cas général. Cependant l'exemple
présenté en (1.5.12) suggère que la détermination des structures de
$\mathbb Z$-catégorie sur une catégorie $\mathcal C$ est
particulièrement simple. C'est en effet le cas. Désignons par
$E\text{-}\operatorname{Cat}$ la $U$-catégorie dont les objets sont les
catégories appartenant à $U$ munies d'une équivalence :

\VerdierSemanticAnchor{ch1-1-5-13-eq1}
\begin{equation}
  T\colon\mathcal C\longrightarrow\mathcal C
  \tag{1.5.13.1}
\end{equation}
et dont les morphismes sont les foncteurs
$F\colon\mathcal C\to\mathcal C'$ muni d'un isomorphisme :

\VerdierSemanticAnchor{ch1-1-5-13-eq2}
\begin{equation}
  m\colon FT\longrightarrow T'F.
  \tag{1.5.13.2}
\end{equation}

On a alors évidemment un foncteur de la $U$-catégorie
$\mathbb Z\text{-}\operatorname{Cat}$ dans la $U$-catégorie
$E\text{-}\operatorname{Cat}$ :

\VerdierSemanticAnchor{ch1-1-5-13-eq3}
\begin{equation}
  R\colon \mathbb Z\text{-}\operatorname{Cat}
  \longrightarrow E\text{-}\operatorname{Cat}.
  \tag{1.5.13.3}
\end{equation}

\VerdierSemanticAnchor{ch1-1-5-13-action}
Le foncteur $R$ associe à une $\mathbb Z$-catégorie
$(\mathcal C,\mathbb Z,T,c)$ la catégorie $\mathcal C$ munie de
l'équivalence $T(1)\colon\mathcal C\to\mathcal C$ et à tout
$\mathbb Z$-foncteur :
\[
  (F,m)\colon(\mathcal C,\mathbb Z,T,c)
  \longrightarrow(\mathcal C',\mathbb Z,T',c'),
\]
le foncteur $F$ muni de l'isomorphisme
$m(1)\colon FT(1)\to T'(1)F$.

\VerdierSemanticAnchor{ch1-1-5-14}
\VerdierSemanticAnchor{ch1-1-5-14-proposition}
\noindent\textbf{Proposition 1.5.14.}\quad
\emph{Le foncteur $R$ est une équivalence de catégories.}

\VerdierSemanticAnchor{ch1-1-5-14-proof}
Nous nous bornerons à donner des indications sur la démonstration. La
proposition (1.5.2) et la remarque (1.5.8) montrent que $R$ est
pleinement fidèle. Il suffit donc de montrer que $R$ est essentiellement
surjectif. Soit donc
$(\mathcal C,T\colon\mathcal C\to\mathcal C)$ une catégorie munie d'une
équivalence. Soient $T^{-1}\colon\mathcal C\to\mathcal C$ un foncteur
quasi-inverse de $T$ et :

\VerdierSemanticAnchor{ch1-1-5-14-eq1}
\begin{equation}
  \begin{aligned}
    u&\colon\operatorname{id}_{\mathcal C}\longrightarrow TT^{-1},\\
    v&\colon\operatorname{id}_{\mathcal C}\longrightarrow T^{-1}T
  \end{aligned}
  \tag{1.5.14.1}
\end{equation}
les isomorphismes d'adjonction.

Les foncteurs $T^n$ et $(T^{-1})^n$ sont, pour tout entier $n>0$,
adjoints l'un de l'autre, d'où des isomorphismes d'adjonction :

\VerdierSemanticAnchor{ch1-1-5-14-eq2}
\begin{equation}
  \begin{aligned}
    u^{(n)}&\colon\operatorname{id}_{\mathcal C}
      \longrightarrow T^n(T^{-1})^n,\\
    v^{(n)}&\colon\operatorname{id}_{\mathcal C}
      \longrightarrow (T^{-1})^nT^n.
  \end{aligned}
  \tag{1.5.14.2}
\end{equation}
